\section*{序論}
11月3日朝から4日朝にかけての生成AI関連X投稿では、AIの用途がソフトウェア開発から機械・電気設計などの物理世界へ広がる動きと、生成結果の真正性を暗号学的に検証する研究が注目された。OpenAIはGPT-6 Astraを活用した試作デバイスの設計を紹介し、Anthropicはモデル推論の証明技術を公開したと報告している。

同時に、AIデータセンターの需要、企業内の発言規制、金融サービスによるAIエージェントのアクセス制限など、AIの社会実装に伴う制度・運用面の課題も話題となった。新モデルの噂や特化型モデルの性能情報には、確認の程度が異なる投稿が混在する。

以下ではGrokが収集した8件を記載順に記事化する。数値や企業発表はGrokの収集内容に基づき、未確認の内容を確定事実として扱わない。独立したWeb検証は実施していない。
